\chapter{Glossary}\label{app:glossary}

This glossary provides quick definitions of the acronyms and recurring terms used in the thesis.
Formal cryptographic definitions appear in \cref{app:definitions}.

\begin{description}
\item[ACI] Account identifier: Signal's service identifier for an account. Before DUET uses an ACI
  as the VRF input or in dual-index and CEA label derivation, the integration converts its string
  form to 16 raw UUID bytes.
\item[ACKA] Authenticated Continuous Key Agreement, the framework of Dowling and Hale whose abstract
  log component this work instantiates and measures (\cref{def:cka}).
\item[AKD] Auditable Key Directory, Meta's deployed key-transparency design for
  WhatsApp~\cite{whatsapp2023whitepaper}.
\item[AKV] Automatic Key Verification, Signal's deployed key-transparency lookup check. It verifies a
  fetched key against the log but carries no peer-warning channel.
\item[Auditor] A verifier that checks signed heads and replays each epoch's declared update set
  against the committed root (the transition audit). Fork detection additionally requires comparison
  with a signed head obtained from an independent vantage
  (\cref{def:epoch-delta,alg:auditor}).
\item[Auditor copy] The auditor's independently reconstructed copy of the committed directory,
  rebuilt by replaying each epoch's declared update set from epoch~0. The transition audit checks
  each committed root against it instead of trusting server-supplied membership proofs
  (\cref{alg:auditor}).
\item[Candidate copy] A temporary clone of the accepted directory state onto which one epoch range is
  replayed. It is committed back only if the whole head passes every check, so a
  rejected epoch leaves the accepted directory state untouched (\cref{alg:auditor}).
\item[CEA / CEAsec] Continuous entity authentication and its security notion within
  ACKA. The continuous operating mode delivers detection in this style.
\item[CKA] Continuous Key Agreement, the abstraction of a ratcheting protocol as a stream of evolving
  keys, used by ACKA.
\item[CKV] Contact Key Verification, Apple's key-transparency mechanism for
  iMessage~\cite{apple-ckv}.
\item[CONIKS] An early key-transparency system that introduced private indices derived through a
  verifiable unpredictable function and a Merkle prefix directory. It is the closest published design point for proof size.
\item[Consistency proof] A short proof that a later log is an append-only extension of an earlier one
  (\cref{def:consistency}).
\item[Declared update set (epoch delta)] The list of inserts and value-disclosing rotations declared
  for a committed epoch and provided for auditor replay. The format cannot express a deletion
  (\cref{def:epoch-delta}).
\item[Deletion evidence] A bundle showing that the operator signed an earlier state containing a
  warning and a later state in the same append-only history in which the same warning slot is absent
  (\cref{def:deletion-evidence}).
\item[Dolev--Yao] The standard active-network adversary model: full control of the wire (intercept,
  modify, inject, delete, replay).
\item[Double Ratchet] Signal's per-message key-ratcheting algorithm (\cref{def:doubleratchet}).
\item[Dual-index warning scheme] This work's construction: two lookup labels and one shared warning
  label per conversation, derived from the shared secret, with the shared warning label carrying an
  authenticated notification after a substitution is detected (\cref{sec:def-dualindex}).
\item[ECVRF] Elliptic-Curve Verifiable Random Function. The suite used is
  ECVRF-EDWARDS25519-SHA512-TAI (\cref{def:vrf}).
\item[Ed25519] The EdDSA signature scheme used to sign the signed tree head (\cref{def:sig}).
\item[ELEKTRA] A multi-device key-transparency system with device cross-signing.
\item[Epoch] A committed version of the directory. A commit appends the prefix root to the log and
  publishes a signed tree head (\cref{def:epoch}).
\item[Equivocation / fork] A malicious server signs two divergent views at one epoch. Detection
  requires comparing signed heads obtained from independent vantage points (A4, under G4 Fork
  Detection).
\item[HKDF] HMAC-based key-derivation function (RFC~5869), used for the dual-index label derivation
  (\cref{def:hkdf}).
\item[Inclusion proof] A proof that a value is present at a position in the log or prefix tree
  (\cref{def:inclusion,def:noninclusion}).
\item[KT] Key transparency: making a key directory's contents publicly and verifiably consistent.
\item[MAC] Message authentication code: a shared-key tag that authenticates a message and detects
  modification. DUET uses HMAC-SHA256 to authenticate warning payloads
  (\cref{def:mac,def:dualindex-write}).
\item[Merkle log] An append-only tamper-evident log whose root commits to every entry in order
  (\cref{def:merkle-log}).
\item[MitM] Man-in-the-middle: an active adversary substituting keys in transit (attack A3).
\item[MLS] Messaging Layer Security, the standardized group-messaging protocol. Group support is
  future work.
\item[Non-inclusion proof] A proof that a label is absent, its slot provably empty
  (\cref{def:noninclusion}). Non-inclusion is a clean empty result on its own. Together with a key
  mismatch, or after prior verified inclusion, it contributes to censorship or removal detection.
\item[OPRF] Oblivious pseudorandom function (RFC~9497). Proposed as future work for query privacy
  from the operator.
\item[Parakeet] A witness-quorum key-transparency system optimized for large-scale deployment.
\item[PNI] Phone-number identifier: Signal's phone-number-scoped service identifier. The current DUET
  integration excludes PNI inputs at the integration point, so they do not enter registration,
  verification, background monitoring, or CEA processing.
\item[PQXDH] The post-quantum variant of X3DH, augmenting the handshake with a key-encapsulation mechanism (KEM).
\item[Prefix tree] A sparse Merkle tree over the space of 256-bit labels, storing the current
  directory (\cref{def:prefix-tree}).
\item[PRF] Pseudorandom function: a keyed function whose outputs are indistinguishable from random
  without knowledge of the key. HMAC-SHA256 serves as the PRF in the continuous-mode state update.
\item[Safety Number] Signal's manual verification code: a 60-digit fingerprint of the two parties'
  identity keys and account identifiers that users are expected to compare out of band. DUET's
  dual-index warning is intended to reduce or replace this ceremony for honestly established and
  monitored conversations.
\item[SEEMless] A key-transparency system built on an append-only zero-knowledge set (aZKS). It
  provides multi-path lookup proofs.
\item[STH] Signed Tree Head: the operator's signed commitment to the log root at an epoch
  (\cref{def:sth}).
\item[STRIDE] A threat-elicitation framework (Spoofing, Tampering, Repudiation, Information
  disclosure, Denial of service, Elevation of privilege). The coverage tables are in
  \cref{sec:sec-stride}.
\item[TOFU] Trust-on-first-use: the initial key is accepted on faith, a limitation inherited from the
  underlying protocol.
\item[Transition audit] The auditor's per-epoch replay of the declared update set against the
  previous committed root, which catches any undeclared mutation, in particular a deletion
  (\cref{def:epoch-delta}, \cref{alg:auditor}).
\item[UNF-CMA] Unforgeability under chosen-message attack: an adversary cannot produce a new valid
  message--tag or message--signature pair, even after obtaining tags or signatures for messages of
  its choice (\cref{def:mac,def:sig}).
\item[VKD] Verifiable Key Directory, the general term for a key directory serving cryptographic proofs
  of its answers.
\item[VRF] Verifiable Random Function: a keyed function whose outputs are pseudorandom to outsiders
  yet provably correct (\cref{def:vrf,def:vrf-security}).
\item[Windowed consistency proof] A single consistency proof spanning an entire offline window.
  Its size is logarithmic in the endpoint log sizes, and it does not require one proof per missed
  epoch (\cref{def:windowed} and \hyperref[lem:windowed-size]{Lemma~\ref*{lem:windowed-size}}).
\item[Write-once policy] The directory's raw-label interface rejects a write to an occupied
  dual-index, CEA, or warning label, so a committed value cannot be overwritten through the honest
  interface (\cref{def:epoch-delta}, \cref{ssec:design-detection}).
\item[X3DH] Extended Triple Diffie--Hellman, Signal's asynchronous key-agreement handshake
  (\cref{def:x3dh}). Its shared secret is the entropy source for the dual-index labels.
\end{description}
